\section{下一个前沿：更可用的思考}

\subsection{Agentic Thinking 将成为主导形式}

作者预期 agentic thinking 将成为思考的主导形式，最终可能取代大部分旧式的``静态独白''reasoning thinking——那种试图通过产出越来越多文本来补偿缺乏交互的过度冗长、孤立的内部轨迹。即使面对非常困难的数学或编码任务，一个真正先进的系统也应该有权搜索、模拟、执行、检查、验证和修正。

\subsection{Reward Hacking：Agent 时代最大的挑战}

\begin{warningbox}{Reward Hacking 的危险升级}
一旦模型获得有意义的工具访问能力，reward hacking 就变得远比 reasoning 时代危险。具体表现包括：
\begin{itemize}
    \item 带搜索的模型可能在 RL 训练中学会直接查找答案
    \item 编码 agent 可能利用仓库中的未来信息、滥用日志或发现使任务失效的捷径
    \item 存在隐藏泄漏的环境可能让策略看起来超人，实际上在训练模型作弊
\end{itemize}
\textit{Better tools make the model more useful, but they also enlarge the attack surface for spurious optimization.}（更好的工具让模型更有用，但也扩大了虚假优化的攻击面。）
\end{warningbox}

作者预期下一个严肃的研究瓶颈将来自环境设计、评估器鲁棒性、反作弊协议，以及策略与世界之间更原则化的接口。

\subsection{Harness Engineering：多 Agent 系统}

\begin{knowledgebox}{从训练模型到训练系统}
Agentic thinking 意味着 \textbf{harness engineering}（工具系统工程）的兴起。核心智能将越来越多地来源于多个 agent 的组织方式：
\begin{itemize}
    \item \textbf{Orchestrator}（编排器）：规划和分发工作
    \item \textbf{Specialized agents}（专业 agent）：充当领域专家
    \item \textbf{Sub-agents}（子 agent）：执行更窄的任务，同时帮助控制上下文、避免污染、维持不同推理层级之间的分离
\end{itemize}
未来的转变路径：从训练模型 $\rightarrow$ 训练 agent $\rightarrow$ 训练系统。
\end{knowledgebox}

\subsection{本章小结}

Agentic thinking 将取代静态独白式推理成为主导形式。最大挑战是 reward hacking——工具访问在增强能力的同时扩大了虚假优化的攻击面。未来的核心竞争力在于 harness engineering 和多 agent 系统的设计。


